% !TEX root = ../../Main.tex
\section{Summary}
\label{Section:Summary}

This work built three things and answered one question.

The first is a trading framework. In it, one strategy implementation runs in three places: an offline backtest over stored quotes, the cTrader backtesting tab, and a live account. Its execution engine reproduces the platform's own results byte for byte on six reference configurations. It also goes further than the platform. It steps through every tick instead of only the bar closes, and it charges the spread recorded at the moment of each trade. The framework reads from a market database of $1.66$ billion quotes for the seven pairs, which takes up $295$ GB. Every bar series used in this work is built from this database.

The second is an agent. It observes thirty-eight dimensionless values, built from sixteen technical indicators. It outputs one continuous action that sets both the direction and the size of its position at the same time. Six changes to the published algorithm were needed to make it train on this problem, and two of them matter most. The first is a scale-sensitive reward. A scale-invariant reward gives no gradient towards holding a position large enough to matter. The second is a penalty on the actor's pre-activation. Without it, the policy saturates at $+1$ or $-1$ and collapses to a constant action.

The third is a selection procedure. First, gates reject candidates with degenerate policies, such as those that hardly trade or keep to one side of the market. Next, each candidate is tested against a null built from itself, by rotating the model's own sequence of positions in time. Last, the candidates are ranked by a composite score in which return is one of nine parts. The procedure was not designed to be elegant. It was built because ranking on return does not work, and \Cref{Section:StatisticalScreening} shows this with measured results.

The answer to the question is as follows. All seven agents are profitable over the eleven-year window, from each of the five starting balances. Their total returns range from $+15.7\%$ to $+226.6\%$ net of the spread and commission of a swap-free raw-spread account. Six of the seven have a positive alpha when their return is decomposed against their own market. Three of these six do this with almost no market exposure. The seventh, USD/JPY, only follows its market. It is reported as a failure, together with a diagnosis of its cause. An equal-weight portfolio of the seven lowers the maximum drawdown from an average of $34.5\%$ to $15.9\%$ and raises the Sharpe ratio to $0.632$. This happens because the agents are close to uncorrelated with one another.

Two findings lie behind these numbers. They are more likely than the numbers themselves to carry over to other problems. First, the agents really use the continuous action. The mean absolute action ranges from $0.09$ to $0.98$ across the pairs, and saturation is below one percent on every pair. So the agents learned to size their positions, and they do not just switch between the two extremes. Second, win rates lie between $49.8\%$ and $55.0\%$, and the strongest model in the set has the lowest win rate. So the profit does not come from being right more often. It comes from what the agents do with a position once it is open. A continuous action makes this possible and a discrete action does not.